\section{Solución Propuesta}
\label{sec:solucion}

La solución que proponemos traduce la pregunta de investigación en un
\emph{pipeline} de cinco etapas encadenadas. Cada etapa resuelve una
dificultad concreta identificada en la sección anterior: el ruido
estadístico que contamina las correlaciones, la densidad excesiva que
uniformiza el espectro, la inestabilidad del orden de los eigenvalores
a lo largo del tiempo y el riesgo de confundir volatilidad agregada de
mercado con manipulación coordinada. El detector no analiza tokens en
aislamiento sino la huella colectiva que deja la coordinación artificial
en la topología de la red.

% -------------------------------------------------------
\subsection{Etapa 1: Recolección y construcción de la matriz de correlación}

Recolectamos velas de 1~minuto para los pares \texttt{SYM/BTC}
involucrados en los ${\approx}351$ eventos confirmados del dataset de
La~Morgia \textit{et al.}~\cite{lamorgia2020}, usando la API pública de
Binance (\texttt{GET /api/v3/klines}) para el período 2021--2022.
Elegimos Binance porque sus datos históricos de este período tienen
cobertura sustancialmente mayor que la de datasets anteriores, lo que
reduce el número de tokens que ya no existen o tienen series
incompletas.

A partir de los precios de cierre calculamos los retornos logarítmicos
$r_i(t) = \ln P_i(t) - \ln P_i(t-1)$ y sobre cada ventana deslizante
de longitud $T$ construimos la matriz de correlación de Pearson
$C_t \in \mathbb{R}^{N \times N}$ según la ecuación~(\ref{eq:pearson}).
Esta es la representación de partida: cada entrada $C_{ij}(t)$ mide
cuánto se mueven juntos los tokens $i$ y $j$ en esa ventana, sin
discriminar todavía qué parte de esa covariación es información real y
qué parte es ruido.

% -------------------------------------------------------
\subsection{Etapa 2: Filtrado espectral con ventana dual}

El problema de usar $C_t$ directamente es que la mayoría de sus entradas
no reflejan relaciones genuinas entre tokens sino fluctuaciones
estadísticas~\cite{plerou1999}. Además, el eigenvalor dominante
concentra la variación conjunta de todos los tokens siguiendo a
Bitcoin ---el modo de mercado global--- y enmascaría cualquier
estructura de comunidad más fina.

Para separar señal de ruido aplicamos la Teoría de Matrices Aleatorias
(RMT). La distribución de Marchenko-Pastur predice que, bajo la hipótesis
de correlaciones puramente aleatorias, todos los eigenvalores caen en el
intervalo $[\lambda_-, \lambda_+]$ definido en la ecuación~\eqref{eq:mp},
con $Q = N/T$ y $\sigma^2$ la varianza promedio de los
retornos~\cite{plerou1999}. Los eigenvalores dentro de ese intervalo son
ruido y se descartan; el eigenvalor más grande $\lambda_1$ (modo de
mercado global) también se elimina.

Aquí surge un problema práctico: la RMT requiere $T \gg N$ para ser
válida, pero un \emph{pump \& dump} dura solo horas~\cite{lamorgia2020},
lo que obliga a usar ventanas de detección cortas donde esa condición
no se cumple. La solución es separar los dos pasos: estimamos los
límites $[\lambda_-, \lambda_+]$ con una ventana larga de calibración
$T_{\text{cal}} = 720~\text{h}$ (30 días, donde la RMT es
estadísticamente válida) y los aplicamos directamente sobre ventanas
cortas de detección de $T_{\text{det}} = 48~\text{h}$. Para verificar
que la aproximación es aceptable, comprobamos que perturbar $\lambda_+$
en ${\pm}15\%$ produce cambios en la precisión del detector menores
al $5\%$.

La matriz filtrada resultante $\tilde{C}$ se obtiene mediante la
ecuación~\eqref{eq:Ctilde}, donde $\mathcal{S} = \{k : \lambda_k >
\lambda_+,\; k \neq 1\}$ retiene únicamente los eigenvalores
informativos, eliminando tanto el ruido aleatorio como el modo de
mercado global.

% -------------------------------------------------------
\subsection{Etapa 3: Construcción del grafo esparsificado}

Con $\tilde{C}$ construimos el grafo de correlaciones usando la distancia
de Mantegna~\cite{mantegna1999} (ecuación~\eqref{eq:mantegna}), con peso
de arista $w_{ij} = 2 - d_{ij}$: tokens más correlacionados reciben
aristas más pesadas. Esta métrica preserva tanto correlaciones positivas
como negativas, a diferencia de truncar los valores negativos a cero.

Usar el grafo completo sobre todos los pares destruye la señal que
queremos detectar: cuando cada nodo está conectado con todos los demás,
los eigenvectores de la Laplaciana tienden a ser uniformes y el IPR
queda cercano a $N^{-1}$ para todos los modos, borrando cualquier
indicio de localización. Para evitarlo reducimos la red en dos pasos.
Primero construimos el Árbol de Expansión Mínima (MST) sobre las
distancias $d_{ij}$, que conecta los $N$ nodos con solo $N-1$ aristas
preservando las correlaciones más fuertes. Luego añadimos las $k$
aristas de mayor peso faltantes en cada nodo, con $k$ elegido para
maximizar la modularidad de la red usando el algoritmo de
Leiden~\cite{traag2019}, que garantiza comunidades bien conectadas y
evita los problemas de comunidades desconectadas del algoritmo de
Louvain. La elección de Leiden es importante porque la optimización de
modularidad es NP-duro en general: Leiden obtiene soluciones de alta
calidad en tiempo polinomial con resultados reproducibles dada una
semilla fija. Verificamos que la estructura espectral se conserva
comparando los primeros 20 eigenvalores antes y después de la reducción
mediante correlación de Spearman ($\rho > 0.9$).

% -------------------------------------------------------
\subsection{Etapa 4: Perfil espectral del IPR robusto al \emph{eigenvalue crossing}}

Sobre el grafo reducido $G_t = (V, E, w)$ calculamos la Laplaciana
normalizada $\mathcal{L}_t$ (ecuación~\eqref{eq:laplacian}). Sus pares
propios $\{\lambda_k, \boldsymbol{\phi}_k\}_{k=0}^{N-1}$ satisfacen
$0 = \lambda_0 \leq \lambda_1 \leq \cdots \leq \lambda_{N-1} \leq 2$.
Para cada eigenvector $k \geq 1$ calculamos el \emph{Inverse Participation
Ratio} (ecuación~\eqref{eq:ipr})~\cite{anderson1958absence}. Un valor
$\mathrm{IPR}_k \approx N^{-1}$ indica participación uniforme de todos
los tokens; un valor cercano a $1$ indica concentración en un solo nodo.
Para un grupo artificial de $m$ tokens el IPR escala como
${\sim}\,m^{-1}$, lo que permite estimar el tamaño del grupo manipulado.

Una forma ingenua de monitorear el IPR sería seguir el eigenvector
número $k$ a lo largo del tiempo. Sin embargo, en redes dinámicas el
orden de los eigenvalores cambia entre ventanas consecutivas: el
eigenvector número~5 en $t$ puede corresponder a una estructura
completamente distinta del número~5 en $t+1$. Comparar
$\mathrm{IPR}_5(t)$ con $\mathrm{IPR}_5(t+1)$ produciría una serie
temporal sin sentido. Este \emph{eigenvalue crossing} es inherente a
perturbaciones transitorias como las que genera un \emph{pump \& dump},
donde la comunidad artificial aparece y desaparece en horas, y no puede
resolverse simplemente restringiendo el análisis a los primeros modos
como se hace en detección de comunidades estáticas.

La solución es no seguir modos individuales. En cada ventana resumimos
la distribución completa del IPR con los estadísticos
$\mathrm{IPR}_{\text{mean}}$ e $\mathrm{IPR}_{p90}$ definidos en las
ecuaciones~\eqref{eq:iprmean}--\eqref{eq:iprp90}. Al no depender del orden de los eigenvalores, ambos son inmunes al
\emph{eigenvalue crossing}. $\mathrm{IPR}_{\text{mean}}$ captura el
nivel promedio de concentración en toda la red; $\mathrm{IPR}_{p90}$
se fija en la cola superior y detecta si existe algún modo muy
localizado sin importar cuál sea. La hipótesis operativa
es que cuando un grupo de tokens está siendo acumulado coordinadamente
antes del \emph{pump}, la comunidad artificial que forman hace que al
menos un modo espectral se concentre sobre ellos, elevando
$\mathrm{IPR}_{p90}(t)$ por encima de lo que cabría esperar de una red
sin manipulación.

% -------------------------------------------------------
\subsection{Etapa 5: Umbral dinámico mediante el \emph{Configuration Model}}

El último problema es definir cuándo un valor de $\mathrm{IPR}_{p90}(t)$
es realmente anómalo. Usar un percentil histórico como umbral es
insuficiente: un evento externo que agite todo el mercado ---como el
colapso de FTX en noviembre de 2022--- puede elevar el IPR en toda la
red sin que haya manipulación coordinada. Un umbral fijo capturaría ese
episodio como una detección, generando exactamente el tipo de falso
positivo que queremos evitar.

Para controlar este efecto usamos el \emph{Configuration
Model}~\cite{newman2003} como punto de comparación ventana a ventana.
Para cada $G_t$ generamos $M = 10\,000$ redes aleatorias con la misma
secuencia de grados pero con las aristas reasignadas uniformemente al
azar, preservando la densidad local de la red pero destruyendo cualquier
estructura de comunidad. Sobre cada red aleatoria calculamos $\mathrm{IPR}_{p90}$ y construimos
la distribución nula empírica. Declaramos una anomalía cuando
$\mathrm{IPR}_{p90}(t)$ supera el umbral definido en la
ecuación~\eqref{eq:threshold}: el percentil~99.7 empírico de esa
distribución. Este criterio no paramétrico evita asumir normalidad, ya
que $\mathrm{IPR}_{p90}$ como estadístico de orden puede seguir una
distribución asimétrica. El umbral es \emph{local}: se recalibra en cada
ventana, de modo que un mercado globalmente más correlacionado produce
automáticamente una línea base más alta. Para evaluar robustez
reportaremos resultados bajo el modelo estándar y bajo una variante con
pesos de arista permutados aleatoriamente.

% -------------------------------------------------------
\subsection{Evaluación e interpretabilidad}

El detector se evalúa sobre los ${\approx}351$ eventos confirmados del
dataset~\cite{lamorgia2020} (pares \texttt{SYM/BTC} en Binance,
2021--2022). La ventana de detección se centra en las 48~h previas a la
hora exacta de inicio de cada evento etiquetado.

Dado el severo desbalance de clases inherente a la detección de anomalías
---los eventos de \emph{pump \& dump} representan una fracción pequeña de
las ventanas temporales totales en dos años de datos---, la métrica
principal es el \textbf{área bajo la curva Precisión-Recall (PR-AUC)},
complementada con el \textbf{F1-Score}. Estas métricas son las estándar
para problemas de detección de fraude con clases desbalanceadas porque no
inflan el resultado con los Verdaderos Negativos, a diferencia del
ROC-AUC. Adicionalmente se reporta ROC-AUC para comparabilidad con la
literatura existente.

Como línea base para comparación usamos un detector de umbral sobre
retornos: se declara anomalía si el retorno promedio de los tokens en la
ventana supera $\mu_{\text{hist}} + 3\,\sigma_{\text{hist}}$, donde
$\mu_{\text{hist}}$ y $\sigma_{\text{hist}}$ se estiman sobre los 30
días previos. Este detector es representativo de los enfoques
univariados actuales~\cite{kamps2018} y permite cuantificar el
beneficio de incorporar estructura de red al análisis.

Más allá de las métricas agregadas, inspeccionaremos los eigenvectores
localizados para verificar que los tokens con mayor peso en
$\boldsymbol{\phi}_{k^*}(t)$ ---donde $k^*$ es el modo que maximiza el
IPR en la ventana de alerta--- corresponden efectivamente a los tokens
manipulados según el dataset. Esta verificación es importante por dos
razones. Primero, da interpretabilidad al resultado: el detector no es
una caja negra sino que señala los nodos específicos que generan la
anomalía. Segundo, si el resultado agregado es negativo ---es decir,
si el IPR no logra separar eventos manipulados de períodos normales---,
este análisis permite determinar si la señal existe pero el umbral no
es suficientemente sensible, o si la manipulación en este conjunto de
datos no produce la estructura espectral que la hipótesis predice. En
ambos casos el resultado es informativo y publicable.

\begin{table}[htbp]
  \caption{Hiperparámetros del pipeline}
  \label{tab:hyperparams}
  \centering
  \begin{tabular}{lll}
    \hline
    \textbf{Parámetro} & \textbf{Valor} & \textbf{Justificación} \\
    \hline
    $T_{\text{cal}}$ & 720~h (30 días) & RMT válida con $T \gg N$ \\
    $T_{\text{det}}$ & 48~h & Duración típica de acumulación~\cite{lamorgia2020} \\
    $M$ & 10\,000 redes & Estima el percentil~99.7 con ${\approx}30$ valores en la cola \\
    Umbral & $\hat{F}_{\text{null}}^{-1}(0.997)$ & Percentil~99.7 empírico (no paramétrico) \\
    Sensibilidad RMT & ${\pm}15\%$ en $\lambda_+$ & Cambio en precisión $< 5\%$ \\
    \hline
  \end{tabular}
\end{table}
